%!TEX encoding = IsoLatin

%%%%%%%%%%%%%%%%%%%%%%%%%%%%%%%%%%%%%%%%%%%%%%%%%%%%%%%%%%%%%%%
%                                                             %
% MINIATURAS MATEM\'ATICAS DE LA GACETA DE LA RSME            %
%                                                             %
%%%%%%%%%%%%%%%%%%%%%%%%%%%%%%%%%%%%%%%%%%%%%%%%%%%%%%%%%%%%%%%

\documentclass[autocontact]{gaceta}

%---------------------
% La Gaceta se har\'a con esto activado, 
% pero no hace falta ahora si no quieres:
\usepackage[T1]{fontenc}
\usepackage{lmodern}

% Puedes usar esto (u otra codificaci\'on) si quieres:
\usepackage[latin1]{inputenc}
\usepackage[spanish]{babel}
%\decimalpoint
%---------------------

%---------------------
% Esto ya lo carga el estilo de La Gaceta:
%\usepackage{amsmath, amsthm, amssymb}
%\usepackage{url} % <-- Para p\'aginas web o similar: \url{...}
%\usepackage{graphicx}
%---------------------

%---------------------
% Carga esto para gr\'aficos si lo necesitas:
% \usepackage{wrapfig}
%---------------------

%---------------------
% Carga esto para gr\'aficos si lo necesitas:
%\usepackage{tikz}
%\usepackage{pgfplots}
%---------------------

%---------------------
% <<< ESTO SE AJUSTAR\'A AL EDITAR CADA N\'UMERO DE LA GACETA >>>
%\provisional
\newcounter{pageforminiatura}
\setcounter{pageforminiatura}{134} % Siempre ser\'a p\'agina par
\setcounter{page}{1} % No se usa (dejarlo como 1)
\journame{La Gaceta de la RSME}
\yearofpublication{2020}
\volume{23}
\issuenumber{1}
%---------------------

%---------------------
% Deja esto as\'{\i}:
%\belongstopart{Art\'{\i}culos} 
%\belongstopart{Secciones} 
%\belongstopart{Noticias de la Sociedad} 
%\belongstopart{Actualidad} 
\nopart
%---------------------

%---------------------
%% Usa este tipo de definiciones si las necesitas:
%\theoremstyle{plain}
%%\newtheorem*{Prop}{Proposici\'on}
%\newtheorem*{Teo}{Teorema}
%\newtheorem*{Lema}{Lema}
%\newtheorem*{Cor}{Corolario}
%---------------------


%---------------------
% Deja esto as\'{\i}: 
\supertitle{\large Miniaturas matem\'aticas de \thejourname\\[1mm]
\hfill\footnotesize
Vol.~\thevolume\ (\theyearofpublication), 
N\'um.~\theissuenumber, P\'ag.~\thepageforminiatura\theprovisional}
%
\noeditor
%---------------------

%---------------------
\title{En recuerdo de Srinivasa Ramanujan (1887--1920).\\
�De cu�ntas formas podemos sumar 176 votos?}
%\author{} % Poner \author{} o no poner nada, da igual.
% Para ahorrar espacio, el autor solo se pone en la firma:
\contact{Juan Luis Varona, Dpto. de Matem�ticas y Computaci�n,
Universidad de La Rioja}{jvarona@unirioja.es}{}
% % http://www.unirioja.es/cu/jvarona/
%---------------------


%%%%%%%%%%%%%%%%%%%%%%%%%%%%%%%%%%%%%%%%%%%%%%%%%%%%%%%%%%%%%%%%%%%%%%%
\begin{document}
%%%%%%%%%%%%%%%%%%%%%%%%%%%%%%%%%%%%%%%%%%%%%%%%%%%%%%%%%%%%%%%%%%%%%%%

%%%%%%%%%%%%%%%%%%%%%%%%%%%%%%%%%%%%%%%%%%%%%%%%%%%%%%%%%%%%%%%%%%%%%%%
% Deja esto as\'{\i}: 
\previousvspace{-1.8cm}
\makesupertitle
\previousvspace{-0.25cm}
\maketitle
\vspace*{-0.4cm}
\thispagestyle{empty} % <<<<< Importante
%%%%%%%%%%%%%%%%%%%%%%%%%%%%%%%%%%%%%%%%%%%%%%%%%%%%%%%%%%%%%%%%%%%%%%%

%% A PARTIR DE AQU\'I HAY LIBERTAD PARA EL AUTOR:

Dado un entero positivo $n$, una partici�n de $n$ es una descomposici�n en enteros positivos cuya suma sea~$n$.
Teniendo en cuenta que dos particiones que solo difieren en el orden de los sumandos no se consideran distintas,
el n�mero de particiones de $n$ se denota $p(n)$. 
Calcular $p(n)$ es un problema de combinatoria nada sencillo si $n$ es grande; �c�mo conseguirlo?
Lo plante� Leibniz en una carta fechada en 1696. Euler (alrededor de 1750) prob� que, para $|x| < 1$
(y tomando $p(0)=1$), se puede escribir
\[
  \prod_{m=1}^{\infty} \frac{1}{1-x^{m}} = \sum_{n=0}^{\infty} p(n) x^n.
\]

%Por otra parte, sean 
%\[
%  g_k = {k(3k-1)}/{2}, \quad k = 0, 1, {-1}, 2, {-2}, 3, {-3},\dots
%\]
%los denominados n�meros pentagonales generalizados (cuando $k \ge 1$, 
%esto responde a cierta distribuci�n de puntos en 
%la superposici�n de $k$ pent�gonos regulares superpuestos, y de ah� el nombre).
%Con esta notaci�n, y si tomamos $p(n) = 0$ para $n<0$, Euler tambi�n prob� que
%%\[
%%  p(n) - p(n-1) - p(n-2) + p(n-5) + p(n-7) + \cdots = 0,
%%\]
%%o bien,
%\[
%  p(n) = \sum_{k=1}^{\infty} (-1)^{k+1} \big(p(n-g_k)+p(n-g_{-k})\big),
% \qquad n \ge 1
%\]
%(son sumas finitas). Esto da una f�rmula recursiva para calcular $p(n)$.

Si adem�s tomamos $p(n)=0$ para $n<0$, Euler tambi�n demostr� que
%\[
%  p(n) - p(n-1) - p(n-2) + p(n-5) + p(n-7) + \cdots = 0,
%\]
%o bien,
\[
  p(n) = \sum_{k=1}^{\infty} (-1)^{k+1} \big(p(n-g_k)+p(n-g_{-k})\big),
 \quad n \ge 1, \quad \text{con }g_k = {k(3k-1)}/{2}
\]
(son sumas finitas). Esto da una f�rmula recursiva para calcular $p(n)$.

Con ella, MacMahon calcul� $p(176) = 476\,715\,857\,290$, que es el n�mero de formas 
de conseguir mayor�a absoluta en el Congreso espa�ol 
(sin restricciones sobre el n�mero de partidos ni sobre los diputados de cada uno);
de hecho, lleg� hasta $p(200)$.

%Us�ndola, MacMahon fue capaz de calcular $p(176) = 476\,715\,857\,290$, que es el n�mero de formas 
%distintas de conseguir mayor�a absoluta en el Congreso de los Diputados espa�ol; 
%% (sin restricciones sobre el n�mero de partidos representados ni sobre los diputados de cada uno)
%de hecho, lleg� hasta $p(200)$.

En el que posiblemente sea su resultado m�s conocido, Hardy y Ramanujan probaron, en 1917, que
\[
p(n) = \frac{1}{2\pi\sqrt{2}} \sum_{k \le \alpha\sqrt{n}} A_{k}(n) \sqrt{k} \,
\frac{d}{dn} \Bigg\{\frac{\exp\Big(\frac{\pi}{k} \sqrt{\frac{2}{3}(n-1/24)}\,\Big)}{\sqrt{n-1/24}}\Bigg\}
 + O\big(n^{-1/4}\big),
\]
con $\alpha>0$ constante arbitraria y 
%(con expresiones m�s compactas que las originales)
(usando expresiones equivalentes a las originales)
\[
% NOTACI�N: A veces (menos) ponen 1 \le h \le k, pero da igual, lo mismo da 0 que k.
% Lo que pasa es es si se pone 0 hay que usar el convenio mcd(0,k) = k.
A_{k}(n) = \!\! \sum_{\begin{smallmatrix}0 \le h < k\\ \operatorname{mcd}(h,k)=1\end{smallmatrix}} 
\!\!\omega_{h,k} \exp\bigg({-2\pi} i\frac{h}{k} n\bigg),
\quad
\omega_{h,k} = \exp\bigg( \pi i \sum_{\mu=1}^{k-1} \frac{\mu}{k} 
\Big( \frac{h\mu}{k} - \Big\lfloor \frac{h\mu}{k} \Big\rfloor  - \frac{1}{2} \Big) \bigg).
\]
%(se usa $\lfloor\,\cdot\,\rfloor$ para denotar la parte entera). 
%Pero eso es una f�rmula asint�tica;
%no es una serie convergente si hacemos $\sum_{k=1}^{\infty}$.
%
Esto aproximaba $p(n)$ incluso mejor de lo previsto,
pero es una f�rmula asint�tica, no una serie convergente si hacemos $\sum_{k=1}^{\infty}$.
Modificando ligeramente el m�todo, Rademacher (1937) dio una representaci�n de $p(n)$ 
por medio de una serie convergente:
\[
p(n) = \frac{1}{\pi\sqrt{2}} \sum_{k=1}^{\infty} A_{k}(n) \sqrt{k} \,
\frac{d}{dn} \Bigg\{\frac{\senh\Big(\frac{\pi}{k} \sqrt{\frac{2}{3}(n-1/24)}\,\Big)}{\sqrt{n-1/24}}\Bigg\},
\]
que resuelve el problema iniciado por Leibniz y Euler de forma muy satisfactoria.

\enlargethispage{0.4cm}

%%%%%%%%%%%%%%%%%%%%%%%%%%% BIBLIOGRAPHY %%%%%%%%%%%%%%%%%%%%%%%%%%%

% Truco para que el t\'{\i}tulo ``Referencias'' no quede tan separado
% ni destaque tanto el tama\~no de letra:
%\let\section\subsection
% (se puede hacer pues no hay secciones ni subsecciones;
% tambi\'en lo suelo usar en la secci\'on de "Problemas y Soluciones")

%\begin{thebibliography}{0}

%% http://ramanujan.sirinudi.org/Volumes/published/ram31.html
%\bibitem{HR-CR1917} 
% G. H. Hardy y S. Ramanujan,
% Une formule asymptotique pour le nombre des partitions de $n$,
% \textit{Comptes Rendus}, 2 Jan. 1917.

%% http://ramanujan.sirinudi.org/Volumes/published/ram36.html
%\bibitem{HR-PLMS1918} 
% G. H. Hardy y S. Ramanujan,
% Asymptotic formul\ae\ in combinatory analysis,
% \textit{Proc. London Mathematical Society}, 2, XVII, 1918, 75--115.
%% Proceedings of the London Mathematical Society

%\bibitem{Rad-1937} 
% Hans Rademacher,
% A convergent series for the partition function $p(n)$,
% \textit{Proc. Nac. Acad. Sci.} 23 (February 1, 1937), no.~2, 78--84. 
%% https://doi.org/10.1073/pnas.23.2.78

%\end{thebibliography}

%%%%%%%%%%%%%%%%%%%%%%%%%%%%%%%%%%%%%%%%%%%%%%%%%%%%%%%%%%%%%%%%%%%%%%%
\end{document}
%%%%%%%%%%%%%%%%%%%%%%%%%%%%%%%%%%%%%%%%%%%%%%%%%%%%%%%%%%%%%%%%%%%%%%%
